% !TeX root=../main.tex
\chapter{بحث و نتیجه‌گیری}
%\thispagestyle{empty}
\label{chap:discussion}

\section{مقدمه}

این نوشتار به یک شکافِ ساختاری در ادبیاتِ مدیریتِ سبدِ پروژه
پرداخت: نبودِ چارچوبی رسمی و از نظرِ محاسباتی قابل‌حل برایِ
تخصیصِ عملیاتیِ بودجه در میانِ سبدی متعهدشده از پروژه‌هایِ
تحتِ قراردادِ \lr{FIDIC} تحتِ عدم‌قطعیتِ جریانِ نقدی. در فصلِ
دوم، شکاف‌هایِ موجود در ادبیاتِ مدیریتِ سبدِ پروژه، بودجه‌بندیِ
جریانِ نقدی، و یادگیریِ تقویتی برایِ تخصیصِ ترتیبیِ منابع
مشخص شد؛ در فصلِ سوم، مسئله در قالبِ یک فرآیندِ تصمیمِ مارکوف
فرمول‌بندی و یک محیطِ سازگار با \lr{Gymnasium} طراحی شد؛ و در
فصلِ چهارم، سازوکارهایِ این محیط از طریقِ چهار پیکربندیِ محک و
مقایسه با بهینه‌یِ اطلاعاتِ‌کامل اعتبارسنجی شد. در این فصل،
یافته‌هایِ این مراحل جمع‌بندی می‌شوند، نوآوریِ تئوریک و عملیِ کار
تشریح می‌گردد، محدودیت‌هایِ تحقیق بیان می‌شود و پیشنهادهایی برایِ
ادامه‌یِ این خط از پژوهش ارائه می‌گردد.

\section{محتوا}

\subsection{جمع‌بندی}
\label{subsec:summary}

سه دست‌آوردِ اصلی از این تحقیق به‌دست آمد.

نخست، محیطِ \code{PortfolioBudgetingEnv} به‌عنوانِ یک محیطِ
متن‌باز و کاملاً پارامتری‌شونده و سازگار با \lr{Gymnasium} طراحی
شد که چرخه‌یِ کاملِ جریانِ نقدیِ \lr{FIDIC} را کدگذاری می‌کند:
پیش‌پرداخت و بازیافتِ تدریجیِ آن، پرداخت‌هایِ نقطه‌یِ عطفِ
وزن‌دهی‌شده بر اساسِ پیشرفت با وزن‌هایِ قابل‌تنظیم و قیودِ
واجد‌شرایط‌بودنِ گام‌زمانی، کسر و آزادسازیِ تضمین، بهره‌یِ
برداشتِ خزانه بر رویِ کسریِ نقدیِ پروژه، و یک دروازه‌یِ فسخِ
دوشرطی با شمارنده‌یِ دوره‌یِ مهلتِ ویژه‌یِ هر پروژه. تنها عنصرِ
تصادفیِ محیط، کاراییِ بودجه‌به‌پیشرفتِ $\eta_{i,t}$ است که در هر
دوره از یک توزیعِ قابل‌تنظیم کشیده می‌شود. تمامیِ پارامترها —
برنامه‌یِ زمانی، قرارداد و سبد — از طریقِ یک واژه‌نامه‌یِ
پیکربندی مشخص می‌شوند که مدخل‌هایِ آن توزیع‌هایِ احتمال هستند و
این امر، تصادفی‌سازیِ حوزه و کالیبراسیون برایِ هر بخشِ
پروژه‌محور را بدونِ تغییرِ کد ممکن می‌سازد. سازگاریِ محیط با
رابطِ برنامه‌نویسیِ \lr{Gymnasium} در ۵۱ آزمون در ۱۰ بخشِ آزمایشی
تأیید شد. دفترِ حسابرسیِ \lr{SQLite} نیز هر تخصیص، رویدادِ
جریانِ نقدی، پرچمِ نقض، و شاخصِ \lr{EVM} را در هر گام ثبت می‌کند
و بدین‌ترتیب رفتارِ هر سیاست به‌طورِ کامل قابلِ تفسیر می‌شود.

دوم، مسئله‌یِ تخصیصِ بودجه‌یِ سبدِ پروژه به‌صورتِ یک فرآیندِ
تصمیمِ مارکوفِ با افقِ متناهیِ $\mathcal{M} = (\mathcal{S},
\mathcal{A}, \mathcal{P}, R, \gamma, T)$ با بُعدِ حالتِ $9N + 1$،
بردارِ مشاهده‌یِ $7N + 2$‌بُعدی که از سیگنال‌هایِ عملی‌شده و نه
شاخص‌هایِ خامِ \lr{EVM} استخراج شده، یک تصویرسازیِ سه‌مرحله‌ایِ
عمل که امکان‌پذیریِ خروجیِ هر سیاستی را تضمین می‌کند، و یک تابعِ
پاداش برابر با جریانِ نقدیِ خالصِ تنزیل‌یافته‌یِ سبد، فرمول‌بندی
شد. تمامیِ انگیزه‌هایِ ساختاری — تکمیل، اجتنابِ از نقض، زمان‌بندیِ
نقطه‌یِ عطف، و بازدارندگیِ از برداشتِ خزانه — بدونِ هیچ جمله‌یِ
جریمه‌یِ کمکی، مستقیماً از سازوکارهایِ مالی سرچشمه می‌گیرند.
ویژگیِ مارکوف دقیقاً برقرار است: حالتِ کاملِ سبدِ پروژه $S_t$ یک
آماره‌یِ کافی برایِ تاریخچه‌یِ قابل‌مشاهده است.

سوم، سه اکتشافیِ حریصانه‌یِ قطعی (\lr{Reach-Plan}، \lr{Catchup-T}،
\lr{Catchup-Next-T}) در برابرِ یک بهینه‌یِ اطلاعاتِ‌کاملِ محاسبه‌شده
به‌صورتِ دستی، در چهار پیکربندیِ محک با پیچیدگیِ ساختاریِ فزاینده و
با دنباله‌هایِ کاراییِ ثابت که تکرارپذیریِ دقیق را تضمین می‌کنند،
ارزیابی شدند. شکافِ بهینگیِ بهترین‌اکتشافی از $3.40$ واحدِ ارزشِ
خالصِ فعلی در اتمِ تک‌پروژه‌ایِ مجزا به $72.15$ واحد در سبدِ
هماهنگ‌شده‌یِ سه‌پروژه‌ای افزایش می‌یابد — افزایشی $21$ برابری. در
پیچیده‌ترین پیکربندی، هر اکتشافیِ حریصانه زیان تولید می‌کند
درحالی‌که بهینه‌یِ اطلاعاتِ‌کامل سودآور باقی می‌ماند؛ این نتیجه
همزمان اعتبارِ سازوکارهایِ محیط و یک هدفِ عملکردیِ مشخص برایِ
عامل‌هایِ یادگرفته‌شده‌یِ آتی را برقرار می‌کند.

\subsection{نوآوری}
\label{subsec:implications}

\paragraph{نوآوریِ تئوری.}
برخلافِ مطالعاتِ پیشین که یا به سطحِ راهبردیِ انتخابِ سبد
می‌پردازند یا بودجه‌بندیِ جریانِ نقدی را مستقل از سازوکارهایِ
قراردادیِ \lr{FIDIC} و بدونِ فرمول‌بندیِ رسمیِ تصمیمِ ترتیبی
بررسی می‌کنند، این تحقیق یک فرمول‌بندیِ بستهِ فرآیندِ تصمیمِ
مارکوف با ویژگیِ مارکوفِ دقیق ارائه می‌دهد که سازوکارهایِ مالیِ
واقعی (پیش‌پرداخت، بازیافت، تضمین، فسخ، بهره‌یِ خزانه) را مستقیماً
در فضایِ حالت، تابعِ پاداش، و تصویرسازیِ عمل کدگذاری می‌کند. محیط
همچنین یک مرجعِ تحلیلیِ قابل‌اثبات — کرانِ بالایِ اطلاعاتِ‌کامل
(گزاره‌یِ \ref{prop:upper_bound}) — فراهم می‌آورد که تا پیش از این
در ادبیاتِ موردبررسی وجود نداشت و امکانِ ارزیابیِ دقیقِ هر سیاست
را بدونِ نیازِ به یک حل‌کننده فراهم می‌کند. از منظرِ پژوهشِ
تصمیم‌گیریِ ترتیبی، این محیط نمونه‌ایِ مسئله‌ای است که همزمان
واقع‌گرا (سازوکارهایِ \lr{FIDIC} کالیبره‌شده با رویه‌یِ تجربیِ
ساخت‌وساز)، رسمی‌شده (فرآیندِ تصمیمِ مارکوفِ بسته‌فرم)، و از نظرِ
محاسباتی قابل‌حل است. غیرقابل‌حل‌بودنِ بهینه‌سازیِ تصادفیِ دقیق —
$K^{NT}$ برگِ درختِ سناریو که برایِ یک نمونه‌یِ سه‌پروژه‌ایِ
متوسط به $K^{72} \approx 10^{50}$ می‌رسد — این محیط را به یک
محکِ خوش‌تعریف برایِ روش‌هایِ یادگیریِ تقویتی و برنامه‌ریزیِ پویایِ
تقریبی تبدیل می‌کند.

\paragraph{نوآوریِ عملی.}
نتایجِ آزمایش نشان می‌دهد که هیچ قاعده‌یِ ثابتی به‌طورِ پیوسته
عملکردِ خوبی ندارد: سیاستِ بهینه به‌طورِ هم‌زمان به مسیرِ کارایی،
ساختارِ پرداختِ نقطه‌یِ عطف، زمان‌بندیِ محرکِ پیش‌پرداخت، و
فازبندیِ پروژه‌هایِ هم‌زمان وابسته است. برایِ یک مدیرِ سبدِ پروژه،
این یافته به این معناست که تخصیصِ مبتنی‌بر قاعده‌یِ تجربی — حتی
قاعده‌ایِ به‌خوبی‌کالیبره‌شده مانندِ \lr{Catchup-T} که کفِ
تلورانسِ قراردادی را دنبال می‌کند — با افزایشِ پیچیدگیِ سبد، ارزشِ
قابل‌توجهی را از دست می‌دهد. توصیه‌یِ عملیِ این تحقیق به صنعت آن
است که تصمیماتِ تخصیصِ بودجه در سبدهایِ چندپروژه‌ای نباید بر اساسِ
قواعدِ تک‌پروژه‌ایِ مستقل گرفته شود، بلکه باید سازوکارِ هماهنگیِ
بودجه‌یِ مشترک و زمان‌بندیِ هم‌زمانِ رویدادهایِ نقدی را صریحاً در
نظر بگیرد. دفترِ حسابرسی نیز یک لایه‌یِ تفسیرپذیریِ آماده‌یِ
استقرار فراهم می‌کند: هر تخصیص، پرداخت، و رویدادِ نقض به‌ازایِ هر
دوره ثبت می‌شود و این محیط را حتی پیش از آموزشِ یک سیاستِ یادگرفته‌شده
نیز به ابزاری برایِ پشتیبانیِ تصمیم تبدیل می‌کند. علاوه‌بر این،
چون هر پارامترِ واژه‌نامه‌یِ پیکربندی می‌تواند به‌صورتِ یک توزیعِ
احتمال مشخص شود و حتیٰ اندازه‌یِ سبد ($N$) خود می‌تواند متغیری
تصادفی و قابل‌تنظیم باشد، این طراحی از تصادفی‌سازیِ حوزه در سطحِ
قسمت \citep{Tobin2017} به‌صورتِ ذاتی پشتیبانی می‌کند؛ سیاستی که بر
رویِ توزیعی از پیکربندی‌ها آموزش ببیند، در هر دوره‌یِ آموزشی با
طیفِ متنوعی از ساختارهایِ پیش‌پرداخت، نیمرخ‌هایِ کارایی، وزن‌هایِ
نقطه‌یِ عطف، و اندازه‌هایِ سبد مواجه می‌شود که پیش‌نیازِ تعمیم‌پذیری
در میانِ انواعِ قراردادها و شرایطِ بازار در عمل است.

\subsection{پیشنهادها}
\label{subsec:future}

محیط و محکِ ارائه‌شده در این تحقیق، چهار محورِ مشخص را برایِ
ادامه‌یِ پژوهش تعریف می‌کنند.

\textbf{آموزشِ عاملِ یادگیریِ تقویتی.}
گامِ بعدیِ طبیعی، آموزشِ یک سیاست — به‌عنوانِ نقطه‌یِ شروع،
\lr{PPO} \citep{Schulman2017} با توجه به پایداریِ آن در فضاهایِ
عملِ پیوسته گزینه‌ایِ مناسب است — بر رویِ توزیعی از پیکربندی‌هایِ
تصادفی‌سازی‌شده و ارزیابیِ آن در برابرِ پیکربندی‌هایِ محکِ بخشِ
ارزیابی است. شکافِ اطلاعاتِ‌کاملِ نشان‌داده‌شده در شکلِ
\ref{fig:gap_plot} هدفِ ارزیابی را مشخص می‌کند؛ سیاستی که بخشِ
قابل‌توجهی از شکافِ \lr{TC-3XMAIF-R} ($72.15$ واحد) را ببندد،
بهبودی معنادار نسبت به تمامیِ معیارهایِ حریصانه به‌شمار می‌رود.

\textbf{سیاست‌هایِ تقویت‌شده با مدل‌هایِ زبانیِ بزرگ.}
بردارِ مشاهده‌یِ $o_t$ ساختاریافته، کم‌بُعد ($7N+2$ مؤلفه) و از نظرِ
معنایی قابلِ‌تفسیر است — هر مؤلفه معنایِ مالی یا قراردادیِ مستقیمی
دارد. این ویژگی آن را به ورودیِ طبیعی برایِ عامل‌هایِ استدلال‌گرِ
مبتنی‌بر مدل‌هایِ زبانیِ بزرگ، یکپارچه‌شده از طریقِ \lr{LangChain}
یا چارچوب‌هایِ مشابه، تبدیل می‌کند، به‌طوری‌که اکتشافی‌هایِ محک
به‌عنوانِ سیگنالِ پاداش برایِ امتیازدهیِ عمل در زمینه عمل کنند.
چالشِ اصلیِ چنین رویکردی، تعمیم‌پذیری در میانِ پیکربندی‌هاست:
سیاستی که بر چهار دنباله‌یِ کاراییِ محک بیش‌برازش شود، به سبدهایِ
جدید منتقل نمی‌شود. طراحیِ پیکربندی‌هایِ تعمیم‌پذیری با تزریقِ
عدم‌قطعیتِ کنترل‌شده — شکل‌هایِ متنوعِ منحنیِ \lr{S}، وزن‌هایِ
ناهمگنِ نقطه‌یِ عطف، ساختارهایِ ترکیبیِ پیش‌پرداخت — گامِ بعدیِ
مناسب برایِ ارزیابیِ استحکامِ سیاست است.

\textbf{کرانِ بالایِ الگوریتمیِ رسمی.}
مقدارِ $Z^*$ که به‌صورتِ دستی محاسبه شده، برایِ چهار پیکربندیِ
کارایی‌ثابت دقیق است اما به تنظیماتِ تصادفی تعمیم نمی‌یابد. یک
تخفیفِ برنامه‌ریزیِ خطی از مسئله‌یِ عرّافِ قطعیِ معادله‌یِ
\ref{eq:ub_obj}، یا یک جوابِ برنامه‌ریزیِ پویا برایِ نمونه‌هایِ
کوچک، کرانی دقیق‌تر و از نظرِ الگوریتمی قابل‌تأیید فراهم می‌کند که
سیاست‌هایِ آموزش‌دیده بتوانند در برابرِ آن مقایسه شوند.

\textbf{استقرار در دنیایِ واقعی.}
یکپارچه‌سازی با \lr{Oracle Primavera P6} یا سامانه‌هایِ اطلاعاتیِ
مدیریتِ سبدِ معادل، امکانِ استقرارِ این محیط را به‌عنوانِ ابزارِ
پشتیبانیِ تصمیم فراهم می‌کند، به‌طوری‌که خروجی‌هایِ زنده‌یِ
\lr{EVM} بردارِ حالت را تأمین کنند و محیط به‌عنوانِ یک شبیه‌سازِ
سناریو برایِ تصمیماتِ تخصیصیِ پیشِ‌رو عمل کند. دفترِ حسابرسیِ
\lr{SQLite} از نظرِ ساختاری با طرحواره‌هایِ استانداردِ گزارش‌دهیِ
\lr{EVM} سازگار است و نیازی به ابزارسازیِ اضافه ندارد.

\subsection{محدودیت‌ها}
\label{subsec:limitations}

فرمول‌بندیِ حاضر سه ساده‌سازیِ عمدی را اعمال می‌کند که هر یک مرزی
طبیعی برایِ گسترشِ آتی به‌شمار می‌رود.

\textbf{استقلالِ پروژه‌ها.} پروژه‌ها تنها از طریقِ بودجه‌یِ مشترکِ
$B_t$ با یکدیگر تعامل دارند. استخرهایِ مشترکِ نیرویِ کار یا
تجهیزات، سرریزهایِ فناورانه، و وابستگی‌هایِ متقابل در سمتِ کارفرما
مدل‌سازی نشده‌اند. این محدودیت خارج از کنترلِ مستقیمِ محقق نیست،
بلکه انتخابی آگاهانه برایِ حفظِ سادگیِ فرمول‌بندی بوده است؛
گسترش به پروژه‌هایِ متصل‌از‌طریقِ‌منابع مستلزمِ افزودنِ موجودیِ
منابع به فضایِ حالت و افزودنِ قیودِ اتصال به تصویرسازیِ عمل خواهد
بود.

\textbf{محدودیت به پرداخت‌هایِ نقطه‌یِ عطف.} ساختارِ پرداخت از
سازوکارِ تأییدِ نقطه‌یِ عطفِ \lr{FIDIC} پیروی می‌کند. قراردادهایِ
مبتنی‌بر زمان، هزینه‌به‌علاوه‌ی سود، و ترکیبی در حالِ حاضر پشتیبانی
نمی‌شوند، هرچند معماریِ واژه‌نامه‌یِ پیکربندی از طریقِ پارامترهایِ
اضافیِ نیمرخِ نقطه‌یِ عطف قابلِ‌گسترش به این حالت‌ها نیز هست.
انتخابِ محدوده‌یِ قراردادیِ موردِمطالعه، منطبق با موضوع و بسترِ
مشخص‌شده‌یِ این تحقیق بوده است.

\textbf{تک‌عنصرِ تصادفی.} کاراییِ بودجه‌به‌پیشرفتِ $\eta_{i,t}$
تمامیِ عدم‌قطعیت را در خود جذب می‌کند. تغییراتِ محدوده‌یِ کار،
شوک‌هایِ قیمتیِ خارجی، و تأخیرهایِ پرداختِ کارفرما به‌صورتِ
برون‌زا در نظر گرفته شده‌اند. این محدودیت، ساختارِ مسئله‌یِ تخصیص
را به‌روشنی مجزا می‌کند؛ گسترش به عدم‌قطعیتِ چندمنبعی مستلزمِ
افزودنِ مؤلفه‌هایِ حالتِ بیشتر برایِ حفظِ ویژگیِ مارکوف خواهد بود.
تأثیرِ این سه محدودیت بر یافته‌ها آن است که مقادیرِ عددیِ گزارش‌شده
در فصلِ چهارم — به‌ویژه شکافِ بهینگی — به سازوکارهایِ مشخص‌شده‌یِ
این محیط وابسته‌اند و تعمیمِ کمّیِ آن‌ها به سایرِ ساختارهایِ
قراردادی یا منابعِ عدم‌قطعیتِ نامدل‌شده، بدونِ گسترشِ متناظرِ محیط
موجه نیست.